% =====================================================================
%  SVM + complexity revision blocks (GLOBECOM 2026 reviewer items
%  P1-#1, P1-#3, P1-#5). Blue = changed relative to the submitted PDF.
%  Paste each block at the location given in its header comment.
% =====================================================================


% ---------------------------------------------------------------------
% BLOCK 0 -- PREAMBLE: add right after  \usepackage{color}
% ---------------------------------------------------------------------
\newcommand{\rev}[1]{\textcolor{blue}{#1}}   % inline revision markup


% ---------------------------------------------------------------------
% BLOCK 1 -- Sec. III-A: replace  \subsection{Preparation for Learning}
%            (the SVM text references Sec.~\ref{subsec:feature}, which
%            is currently undefined and prints "Sec. ??")
% ---------------------------------------------------------------------
\subsection{Preparation for Learning}
\label{subsec:feature}


% ---------------------------------------------------------------------
% BLOCK 2 -- Sec. III-B-1: replace everything from
%            \subsubsection{SVM-based Model for TAS}
%            up to (not including) \subsubsection{NB-based Model for TAS}
% ---------------------------------------------------------------------
\subsubsection{SVM-based Model for TAS}
\label{subsubsec:svm}
{\color{blue}
The SVM model constructs separating hyperplanes in a high-dimensional
feature space, which can be used for classification. Among the available
multi-class SVM strategies, namely one-versus-rest (OVR) and one-versus-one
(OVO), we employ the OVO approach to address the antenna selection
problem in \eqref{eq:ETAS_Defn}. Accordingly, $N_S(N_S-1)/2$ binary
classifiers are trained, where the classifier indexed by the pair
$(i,j)$, $1\leq i<j\leq N_S$, is trained only on the subset
$\mathcal{M}_{ij}\subseteq\{1,\dots,M\}$ of training samples labeled as
class $i$ or class $j$, and discriminates between these two classes.
OVO trains more binary classifiers than OVR, which trains only $N_S$.
However, each OVO classifier uses only $|\mathcal{M}_{ij}|\approx 2M/N_S$
samples when the antenna labels are equiprobable, whereas every OVR
classifier uses all $M$ samples. Since kernel SVM training scales at
least quadratically with the number of training
samples~\cite{Bottou_Lin_2007_LSKM}, the total training cost of OVO is
lower than that of OVR, and the gap widens as $N_S$
grows~\cite{Hsu_Lin_2002_TNN}. For $N_S=2$, both strategies reduce to a
single binary classifier.
Prior to training, each element of the normalized feature vector
$\mathbf{t}_m$ in \eqref{norm_feature_vec} is further standardized to
zero mean and unit variance using the sample statistics of the training
set. With a slight abuse of notation, we retain the symbol $\mathbf{t}_m$
for the standardized vector. The parameters of the $(i,j)$-th classifier
are obtained by solving the following regularized hinge-loss minimization
problem~\cite{He_2018_wcLet,Hsu_Lin_2002_TNN}
\begin{multline}\label{eq:SVM}
\{\mathbf{u}_{ij},\kappa_{ij}\} = \argmin_{\mathbf{u}_{ij},\kappa_{ij}}
C \!\!\sum_{m\in\mathcal{M}_{ij}}\!\! \Big[ w_{ij}[m]\,\text{cost}_1\!\left(\mathbf{u}_{ij}^T \mathbf{f}(\mathbf{t}_m)+\kappa_{ij}\right) \\
 + (1-w_{ij}[m])\,\text{cost}_0\!\left(\mathbf{u}_{ij}^T \mathbf{f}(\mathbf{t}_m)+\kappa_{ij}\right) \Big]
+ \tfrac{1}{2}\mathbf{u}_{ij}^T\mathbf{K}_{ij}\mathbf{u}_{ij},
\end{multline}
where $\mathbf{u}_{ij}\in\mathbb{R}^{|\mathcal{M}_{ij}|\times 1}$ and
$\kappa_{ij}$ are the learning parameter vector and bias of the
$(i,j)$-th classifier, and $C$ is a penalty factor that balances the
training error against the regularization term. The vector
$\mathbf{f}(\mathbf{t}_m)\in\mathbb{R}^{|\mathcal{M}_{ij}|\times 1}$
denotes the kernel feature map of the $m^{\text{th}}$ training sample.
It is obtained with a Gaussian radial basis function (RBF) kernel, and
its $q^{\text{th}}$ element, $q\in\mathcal{M}_{ij}$, is given by
$f_q(\mathbf{t}_m) = \exp\!\left(-\|\mathbf{t}_m -\mathbf{t}_q\|^2/(2\sigma^2)\right)$,
where $\sigma^2$ is the kernel width. The matrix $\mathbf{K}_{ij}$ is the
corresponding kernel (Gram) matrix, whose columns are
$\mathbf{f}(\mathbf{t}_m)$, $m\in\mathcal{M}_{ij}$. The term
$\tfrac{1}{2}\mathbf{u}_{ij}^T\mathbf{K}_{ij}\mathbf{u}_{ij}$ is the
squared norm of the separating hyperplane in the kernel-induced feature
space, so minimizing it maximizes the margin. Further, $w_{ij}[m]=1$ if
the $m^{\text{th}}$ training CSI sample is labeled as class $i$, and
$w_{ij}[m]=0$ if it is labeled as class $j$. The hinge-type losses are
defined as $\text{cost}_k(z)=\max\!\left((-1)^k z+1,\,0\right)$ for
$k\in\{0,1\}$, so that $\text{cost}_1(\cdot)$ and $\text{cost}_0(\cdot)$
penalize margin violations of class $i$ and class $j$, respectively.
% !!! Replace with what the simulation code actually did !!!
The hyperparameters $C$ and $\sigma^2$ are selected by $5$-fold
cross-validation on the training set, which yields
$C=$~[\,\textbf{value}\,] and $\sigma^2=$~[\,\textbf{value}\,].
Problem \eqref{eq:SVM} is the standard soft-margin kernel SVM. It admits
a sparse solution in which $u_{ij,q}\neq 0$ only for the training samples
that lie on or inside the margin, i.e., the support
vectors~\cite{Bottou_Lin_2007_LSKM}. Let
$\mathcal{V}_{ij}\subseteq\mathcal{M}_{ij}$ denote the support-vector set
of the $(i,j)$-th classifier, and let
$N_{\mathrm{SV}}=\big|\bigcup_{i<j}\mathcal{V}_{ij}\big|\leq M$ denote
the number of distinct support vectors. Only these $N_{\mathrm{SV}}$
vectors are stored after training.

After training, the parameters $\{\mathbf{u}_{ij},\kappa_{ij}\}$ are
obtained for all $1\leq i<j\leq N_S$, and the SVM model can then perform
antenna selection. For a new CSI instance, a feature vector
$\mathbf{t}\in \mathbb{R}^{1\times N}$ is generated following the feature
extraction process of Sec.~\ref{subsec:feature} and standardized using
the stored training statistics. The decision value of the $(i,j)$-th
classifier is computed from its support vectors as
\begin{equation}\label{eq:SVM_dec}
z_{ij}(\mathbf{t}) = \sum_{q\in\mathcal{V}_{ij}} u_{ij,q}\,
\exp\!\left(-\frac{\|\mathbf{t}-\mathbf{t}_q\|^2}{2\sigma^2}\right) + \kappa_{ij},
\end{equation}
and each binary classifier casts a vote
\begin{equation}\label{eq:SVM_vote}
v_{ij}(\mathbf{t}) =
\begin{cases}
i, & z_{ij}(\mathbf{t}) \geq 0,\\
j, & \text{otherwise}.
\end{cases}
\end{equation}
The predicted antenna index is determined by the OVO majority-vote rule
\begin{equation}\label{eq:SVM_decision}
l^{\ast} = \argmax_{1\leq l \leq N_S} \ \sum_{1\leq i<j\leq N_S}
\mathbf{1}\!\left[v_{ij}(\mathbf{t}) = l\right],
\end{equation}
where $\mathbf{1}[\cdot]$ denotes the indicator function.
}


% ---------------------------------------------------------------------
% BLOCK 3 -- Sec. V: replace the whole section, from
%            \section{Computational complexity Analysis}
%            up to (not including) \begin{table}[t]
%            (the feedback-overhead sentence is kept as-is: it belongs
%            to reviewer item P1-#2, not to this revision)
% ---------------------------------------------------------------------
\section{Computational Complexity Analysis}
\label{sec:complexity}
We summarize the feedback overhead and selection complexity of different
TAS \rev{schemes} in Table~\ref{tab:selection_comparison_modified}. Here,
$N=N_S(N_D+N_E+2)$ for the active eavesdropping scenario, while
$N=N_S(N_D+2)$ for the passive eavesdropping case. We observe that the
ML-based schemes require only half the feedback overhead compared to the
conventional scheme.
{\color{blue}
We distinguish the online complexity, incurred in every coherence block
to select the antenna, from the offline complexity of generating the
labeled dataset and training the classifiers, which is incurred only once.

\emph{Online selection complexity:} The ETAS rule computes
$\gamma_{S_jD}$, and $\gamma_{S_jE}$ under active eavesdropping, for all
$N_S$ antennas in a single pass over the $N$ channel gains, which costs
$\mathcal{O}(N)$. It then finds the maximum of $N_S$ rates by a linear
search, which costs $\mathcal{O}(N_S)$ and requires no sorting. Its
overall complexity is therefore $\mathcal{O}(N)$.
For the SVM-based scheme, the $N$-dimensional RBF kernel in
\eqref{eq:SVM_dec} is evaluated once for each of the $N_{\mathrm{SV}}$
distinct support vectors and shared by all $N_S(N_S-1)/2$ OVO
classifiers, which costs $\mathcal{O}(N_{\mathrm{SV}}N)$. Forming the
decision values costs $\mathcal{O}(N_{\mathrm{SV}}N_S)$, since each
support vector belongs to at most $N_S-1$ classifiers, and the voting in
\eqref{eq:SVM_decision} costs $\mathcal{O}(N_S^2)$. Since $N>N_S$ and
$N_{\mathrm{SV}}\geq N_S$, the overall complexity is
$\mathcal{O}(N_{\mathrm{SV}}N)$, and the OVR strategy has the same
dominant term. Because $N_{\mathrm{SV}}\leq M$ and typically grows
linearly with $M$~\cite{Bottou_Lin_2007_LSKM}, this complexity is upper
bounded by $\mathcal{O}(MN)$.
% !!! Fill in from the trained model !!!
In our simulations, $N_{\mathrm{SV}}=$~[\,\textbf{value}\,] for $M=10000$.
The NB-based scheme with Gaussian class-conditional densities evaluates
$N$ univariate likelihoods for each of the $N_S$ classes in
\eqref{eq:NB_2} and then finds the maximum by a linear search, which
costs $\mathcal{O}(N_SN)$. This cost does not depend on $M$, since only
the $2N_SN$ class-conditional means and variances are stored.
The KNN-based scheme computes the Euclidean distance from $\mathbf{t}$ to
all $M$ stored training vectors, which costs $\mathcal{O}(MN)$. It then
selects the $K$ nearest neighbors and takes a majority vote, which costs
$\mathcal{O}(M)$. Its overall complexity is therefore $\mathcal{O}(MN)$.

\emph{Offline complexity:} All ML-based schemes require labels for the
$M$ training samples. These are obtained by applying the ETAS rule to
each sample, at a total cost of $\mathcal{O}(MN)$. Training the NB model
requires computing the class-conditional means and variances, which
costs $\mathcal{O}(MN)$. KNN requires no training beyond storing the $M$
feature vectors. For the SVM, each of the $N_S(N_S-1)/2$ OVO problems
involves $|\mathcal{M}_{ij}|\approx 2M/N_S$ samples. Evaluating its
kernel matrix therefore costs $\mathcal{O}(M^2N/N_S^2)$, which amounts to
$\mathcal{O}(M^2N)$ over all classifiers. Under OVR, every binary problem
uses all $M$ samples, so this cost rises to
$\mathcal{O}(N_SM^2N)$, which motivates the OVO strategy adopted in
Section~\ref{subsubsec:svm}.

Table~\ref{tab:selection_comparison_modified} shows that the ETAS rule
has the lowest online complexity. Among the ML-based schemes, NB has the
lowest online complexity, and it is independent of $M$. The complexity of
the SVM-based scheme grows with the number of support vectors, while that
of KNN grows linearly with the training-set size $M$. Hence, the ML-based
schemes do not reduce the online computational complexity relative to
ETAS. Instead, the training-set size and, for the SVM, the number of
support vectors govern a tradeoff between secrecy performance and
selection complexity.
}


% ---------------------------------------------------------------------
% BLOCK 4 -- Table I: replace the whole active \begin{table}[t] ...
%            \end{table} that follows Sec. V (keep the commented-out
%            older tables as they are). Fixes the broken "NsvN" cell,
%            the O(N^2)/O(N) entries, and adds offline cost + storage.
% ---------------------------------------------------------------------
\begin{table}[t]
	\caption{Complexity and Overhead of Different TAS Schemes
	\rev{($M$: training-set size; $N_{\mathrm{SV}}\leq M$: number of support vectors)}}
	\label{tab:selection_comparison_modified}
	\centering
	\resizebox{\columnwidth}{!}{
		\begin{tabular}{|c|c|c|c|c|}
			\hline
			\textbf{Metric} & \textbf{ETAS} & \textbf{SVM \rev{(OVO)}} & \textbf{NB} & \textbf{KNN}\\
			\hline
			\textbf{Feedback Overhead}& $N$ complex values & \multicolumn{3}{c|}{$N$ real values}\\
			\hline
			\makecell{\textbf{\rev{Online} Selection}\\\textbf{Complexity}} &
			\rev{$\mathcal{O}(N)$} & \rev{$\mathcal{O}(N_{\mathrm{SV}}N)$} &
			\rev{$\mathcal{O}(N_SN)$} & \rev{$\mathcal{O}(MN)$} \\
			\hline
			\rev{\makecell{\textbf{Offline Complexity}\\\textbf{(labeling + training)}}} &
			\rev{--} & \rev{$\mathcal{O}(M^2N)$} & \rev{$\mathcal{O}(MN)$} & \rev{$\mathcal{O}(MN)$} \\
			\hline
			\rev{\textbf{Storage}} &
			\rev{--} & \rev{$\mathcal{O}(N_{\mathrm{SV}}N)$} & \rev{$\mathcal{O}(N_SN)$} & \rev{$\mathcal{O}(MN)$} \\
			\hline
	\end{tabular}	}
\vsp{-5mm}
\end{table}


% ---------------------------------------------------------------------
% BLOCK 5 -- Remove the old "comparable or lower complexity" claims,
%            which are false once the complexities are corrected.
%            Only the complexity clause is changed; the feedback-overhead
%            wording is left for whoever handles reviewer item P1-#2.
% ---------------------------------------------------------------------

% 5a -- Abstract, replace the last sentence
%       ("Finally, we show that the ML-based schemes require comparable
%        or lower computational complexity and reduced feedback ...")
Finally, we show that the ML-based schemes require reduced feedback
overhead\rev{, and we quantify their online and offline computational
complexity: NB has the lowest online complexity among the ML-based
schemes, whereas the complexity of SVM and KNN scales with the
training-set size,} thereby illustrating performance versus complexity
tradeoffs \rev{with respect to the ETAS rule}.

% 5b -- Introduction, 3rd contribution bullet, replace the last sentence
%       ("Our analysis show that the proposed ML-based approaches require
%        the half amount of feedback overhead and comparable or lower ...")
Our analysis show that the proposed ML-based approaches require the half
amount of feedback overhead \rev{as compared to the ETAS rule, at the cost
of a higher online selection complexity that grows with the number of
support vectors for SVM and with the training-set size for KNN, whereas
the complexity of NB is independent of the training-set size}.

% 5c -- Conclusion, replace the last sentence
%       ("The proposed ML-based approaches require the half amount of
%        feedback overhead and comparable or lower computational ...")
The proposed ML-based approaches require the half amount of feedback
overhead as compared to exhaustive TAS rule\rev{, at the cost of a higher
online selection complexity, which grows with the training-set size for
SVM and KNN,} while achieving near-optimal secrecy performance under
interference constraints and energy harvesting.

% 5d -- (optional) Sec. II-B "Challenges" paragraph, replace the last
%       sentence ("Machine learning models, which can learn complex
%        relations ... transferring the complexity to the offline training phase.")
Machine learning models, which can learn complex relations in a given
data, can thus be exploited for TAS\rev{, with the computationally
intensive training performed offline}.


% ---------------------------------------------------------------------
% BLOCK 6 -- BibTeX: append to Bibtex/priyanka.bib
% ---------------------------------------------------------------------
% @article{Hsu_Lin_2002_TNN,
%   author  = {C.-W. Hsu and C.-J. Lin},
%   title   = {A comparison of methods for multiclass support vector machines},
%   journal = {IEEE Trans. Neural Netw.},
%   volume  = {13},
%   number  = {2},
%   pages   = {415--425},
%   month   = mar,
%   year    = {2002}
% }
%
% @incollection{Bottou_Lin_2007_LSKM,
%   author    = {L. Bottou and C.-J. Lin},
%   title     = {Support vector machine solvers},
%   booktitle = {Large Scale Kernel Machines},
%   editor    = {L. Bottou and O. Chapelle and D. DeCoste and J. Weston},
%   publisher = {MIT Press},
%   address   = {Cambridge, MA, USA},
%   pages     = {1--27},
%   year      = {2007}
% }
